\subsection{\texorpdfstring{空间 \(\mathcal{F}_{\mathfrak{G}}(\mathbf{X};\mathbf{Y})\) 的性质}{空间 F 下标 G (X;Y) 的性质}}

命题 2. 设 \(\mathbf{X}_1, \mathbf{X}_2\) 是两个集，\(\mathbf{Y}\) 是一致空间，\(\mathfrak{S}_i\) 是 \(\mathbf{X}_i\) 的子集组成的集（\(i = 1, 2\)），而 \(\mathfrak{S}_1 \times \mathfrak{S}_2\) 是 \(\mathbf{X}_1 \times \mathbf{X}_2\) 的形如 \(\mathbf{A}_1 \times \mathbf{A}_2\) 的子集组成的集，其中 \(\mathbf{A}_i \in \mathfrak{S}_i\)（\(i = 1, 2\)）。则典范双映射

\[
\mathscr {F} \left(\mathrm{X} _ {1} \times \mathrm{X} _ {2}; \mathrm{Y}\right)\rightarrow \mathscr {F} \left(\mathrm{X} _ {1}; \mathscr {F} \left(\mathrm{X} _ {2}; \mathrm{Y}\right)\right)
\]

（集合论，R，§ 4，第 14 号）是一致空间

\[
\mathcal {F} _ {\mathfrak {S} _ {1} \times \mathfrak {S} _ {2}} (\mathrm{X} _ {1} \times \mathrm{X} _ {2}; \mathrm{Y})
\]

到 \(\mathcal{F}_{\mathfrak{S}_1}(\mathbf{X}_1;\mathcal{F}_{\mathfrak{S}_2}(\mathbf{X}_2;\mathbf{Y}))\) 上的同构

设 \(\mathbf{V}\) 是 \(\mathbf{Y}\) 的围域，\(\mathbf{A}_i\in \mathfrak{S}_i\)（\(i = 1,2\)）；则由定义立即可见，\(W(\mathbf{A}_1\times \mathbf{A}_2,\mathbf{V})\) 经典范双映射与 \(W(\mathbf{A}_1,W(\mathbf{A}_2,\mathbf{V}))\) 等同，结论立得。

命题 3. a) 设 X 是一个集；S 是 X 的子集组成的一个集；Y、Y' 是两个一致空间；f: Y → Y' 是一致连续映射。则 F\_S(X; Y) 到 F\_S(X; Y') 中的映射 u → f ∘ u 是一致连续的。

\begin{enumerate}
\def\labelenumi{\alph{enumi})}
\setcounter{enumi}{1}
\tightlist
\item
  设 \(\mathbf{X}, \mathbf{X}'\) 是两个集；\(\mathfrak{S}\)（相应地，\(\mathfrak{S}'\)）是 \(\mathbf{X}\)（相应地，\(\mathbf{X}'\)）的子集组成的集；\(\mathbf{Y}\) 是一致空间；\(g: \mathbf{X}' \to \mathbf{X}\) 是这样的映射：对每个 \(\mathbf{A}' \in \mathfrak{S}'\)，\(g(\mathbf{A}')\) 含于 \(\mathfrak{S}\) 中集的有限并。则映射 \(u \to u \circ g\)（\(\mathcal{F}_{\mathfrak{S}}(\mathbf{X}, \mathbf{Y})\) 到 \(\mathcal{F}_{\mathfrak{S}'}(\mathbf{X}', \mathbf{Y})\) 中的）是一致连续的。
\end{enumerate}

命题 4. 设 X、Y 是两个集，\((\mathbf{X}_{\lambda})_{\lambda \in \mathbf{L}}\) 是一个集族，\((\mathrm{Y}_{\mu})_{\mu \in \mathbf{M}}\) 是一个一致空间族。对每个 \(\lambda \in \mathbf{L}\)，设 \(\mathfrak{S}_{\lambda}\) 是 \(\mathbf{X}_{\lambda}\) 的子集组成的集，\(g_{\lambda}\) 是 \(\mathbf{X}_{\lambda}\) 到 X 中的映射，而 \(\mathfrak{S}\) 是 X 的子集组成的集，即诸集 \(g_{\lambda}(\mathfrak{S}_{\lambda})\) 的并。对每个 \(\mu \in \mathbf{M}\)，设 \(f_{\mu}\) 是 Y 到 \(Y_{\mu}\) 中的映射，并赋予 Y 使诸 \(f_{\mu}\) 都一致连续的最粗一致结构。则 \(\mathcal{F}(X;Y)\) 上 S-收敛的一致结构是使映射 \(u \to f_{\mu} \circ u \circ g_{\lambda}\)（\(\mathcal{F}(X;Y)\) 到 \(\mathcal{F}_{\mathfrak{S}_{\lambda}}(X_{\lambda},Y_{\mu})\) 中的）都一致连续的最粗一致结构。

这些命题是第 2 号注记 1 中给出的 S-收敛的一致结构的围域基本系描述的直接推论；证明的细节留给读者。特别地，命题 4 蕴含：

推论. 设 X 是一个集，\((\mathrm{Y}_{t})_{t\in\mathbf{I}}\) 是一致空间族，S 是 X 的子集组成的集。若赋予 \(\prod_{t\in I}Y_{t}\) 乘积一致结构，则一致空间 \(\mathcal{F}_{\mathfrak{S}}\left(\mathrm{X},\prod_{t\in\mathbf{I}}\mathrm{Y}_{t}\right)\) 到乘积一致空间 \(\prod_{t\in\mathbf{I}}\mathcal{F}_{\mathfrak{S}}(\mathrm{X};\mathrm{Y}_{t})\) 上的典范双映射（集合论，R，§ 4，第 13 号）是同构。

